
les morphismes @@M0@@ pour @@M1@@ (resp. @@M2@@) des fibres spéciales et générales dans le schéma ambiant. Nous désignons par
@@M3@@
les morphismes induits, donnant donc lieu aux carrés cartésiens
@@M4@@

Supposons d'abord @@M5@@ \emph{propre}. Je dis qu'alors le diagramme suivant est commutatif :
@@M6@@
Ce résultat généralise 7.5 (car pour @@M7@@, les flèches de (7.3.2) s'identifient à la flèche identique de @@M8@@). On notera que si @@M9@@ et @@M10@@ satisfont aux conditions de 7.7 i.e. sont à fibres régulières, alors @@M11@@ et @@M12@@ sont automatiquement de tor-dimension finie, et le diagramme peut s'interpréter comme un diagramme commutatif avec des @@M13@@, qu'on laisse au lecteur de récrire.

Pour prouver la commutativité de (7.10.1), on voit aussitôt, en utilisant la définition de @@M14@@ et @@M15@@, qu'on est ramené à prouver la commutativité des deux carrés antérieurs du prisme
@@M16@@
